% SEGMENT OLP-0400-S01
% Part: many-valued-logic
% SOURCE CORRECTION TA-MVL-017: correct the chapter metadata from three-valued-logics to infinite-valued-logics.
% Chapter: infinite-valued-logics
% Section: lukasiewicz

\documentclass[../../../include/open-logic-section]{subfiles}

\begin{document}

\olfileid{mvl}{inf}{luk}

\olsection{லூகாசியேவிச்~(\L ukasiewicz) தருக்கம்}

\begin{editorial}
  இது முடிவிலா மதிப்பு லூகாசியேவிச்~(\L ukasiewicz) தருக்கத்தைப் பற்றிய
  சுருக்கமான, முழுமையடையாத பிரிவு.
\end{editorial}

\begin{defn}\ollabel{def:lukasiewicz} முடிவிலா மதிப்பு
லூகாசியேவிச்~(\L ukasiewicz) தருக்கம்~$\LogLuk[\infty]$ பின்வரும்
அணியைக் கொண்டு வரையறுக்கப்படுகிறது:
\begin{enumerate}
  \item $\lnot$, $\land$, $\lor$, $\lif$ ஆகியவற்றைக் கொண்ட செந்தரக்
  கூற்று மொழி~$\Lang L_0$.
  \item மெய்மதிப்புகளின் கணம்~$V_\infty$.
  \item $1$ மட்டுமே நியமிக்கப்பட்ட மதிப்பு; அதாவது, $V^+ = \{1\}$.
  \item வாய்மைச் சார்புகள் பின்வரும் சார்புகளால் தரப்படுகின்றன:
  \begin{align*}
    \tf{\lnot}[\LogLuk](x) & = 1 - x\\
    \tf{\land}[\LogLuk](x,y) & = \min(x,y)\\
    \tf{\lor}[\LogLuk](x,y) & = \max(x,y)\\
    \tf{\lif}[\LogLuk](x,y) & = \min(1,1-(x-y)) = \begin{cases}
      1 & \text{எனில் } x \le y\\
      1-(x-y) & \text{வேறெனில்.}
    \end{cases}
    \end{align*}
\end{enumerate}
$m$-மதிப்பு லூகாசியேவிச்~(\L ukasiewicz) தருக்கமும் இதேவாறு
வரையறுக்கப்படுகிறது; ஆனால் அதில் $V = V_m$.
\end{defn}

\begin{prop}
  \olref[thr][luk]{def:lukasiewicz} இல் வரையறுக்கப்பட்ட தருக்கம்
  $\LogLuk[3]$, \olref{def:lukasiewicz} இல் வரையறுக்கப்பட்ட
  $\LogLuk[3]$ உடன் ஒன்றே.
\end{prop}

\begin{proof}
  \olref[thr][luk]{def:lukasiewicz} இல் இணைப்பிகளுக்காகத் தரப்பட்ட வாய்மை
  அட்டவணைகளையும், \olref{def:lukasiewicz} இலுள்ள சமன்பாடுகள் தீர்மானிக்கும்
  பின்வரும் வாய்மை அட்டவணைகளையும் ஒப்பிட்டால் இதைக் காணலாம்:
  \begin{center}
    \begin{tabular}{c|c} 
      $\tf{\lnot}$ & \\ 
      \hline  
      $1$ & $0$ \\ 
      $1/2$ & $1/2$ \\
      $0$ & $1$ 
    \end{tabular}
    \quad
    \begin{tabular}{c|ccc} 
      $\tf{\land}[\LogLuk[3]]$ & $1$ & $1/2$ & $0$ \\ 
      \hline 
      $1$ & $1$ & $1/2$ & $0$ \\ 
      $1/2$ & $1/2$ & $1/2$ & $0$\\ 
      $0$ & $0$ & $0$ & $0$ 
    \end{tabular}
    \\[2ex]
    \begin{tabular}{c|ccc} 
      $\tf{\lor}[\LogLuk[3]]$ & $1$ & $1/2$ & $0$ \\ 
      \hline 
      $1$ & $1$ & $1$ & $1$ \\ 
      $1/2$ & $1$ & $1/2$ & $1/2$ \\
      $0$ & $1$ & $1/2$ & $0$ 
    \end{tabular}
    \quad
    \begin{tabular}{c|ccc} 
      $\tf{\lif}[\LogLuk[3]]$ & $1$ & $1/2$ & $0$ \\ 
      \hline 
      $1$ & $1$ & $1/2$ & $0$ \\ 
      $1/2$ & $1$ & $1$ & $1/2$  \\ 
      $0$ & $1$ & $1$ & $1$ 
    \end{tabular}
  \end{center} 
\end{proof}

\begin{prop}\ollabel{prop:luk-infty-m}
  $\Gamma \Entails[\LogLuk[\infty]] !B$ எனில், எல்லா~$m \ge 2$ க்கும்
  $\Gamma \Entails[\LogLuk[m]] !B$.
\end{prop}

\begin{proof}
  பயிற்சி.
\end{proof}

\begin{prob}
  \olref[mvl][inf][luk]{prop:luk-infty-m} ஐ நிறுவுக.
\end{prob}

உண்மையில், இதன் மறுதலையும் அமைகிறது.

முடிவிலா மதிப்பு லூகாசியேவிச்~(\L ukasiewicz) தருக்கமே மிகவும் பரவலாகப்
பயன்படுத்தப்படும் மங்கல் தருக்கம். மங்கல் தருக்க இலக்கியத்தில்,
நிபந்தனைவாக்கியம் பெரும்பாலும் $\lnot !A \lor !B$ என வரையறுக்கப்படுகிறது.
அவ்வாறு வரையறுத்தால் முடிவிலா மதிப்பு வலு க்ளீனி தருக்கம் கிடைக்கும்.

\begin{prob}
  $(p \lif q) \lor (q \lif p)$ என்பது~$\LogLuk[\infty]$ இன் ஒரு
  மெய்மம் என்று காட்டுக.
\end{prob}

\end{document}
